% GENERATED FILE. Edit canonical lessons, then run npm run generate:books.
% canonical-source-hash: fnv1a64:493ebe0278d9118b
% canonical-lessons: BN-C307-dhak, BN-C307-kasor, BN-C307-dhunuchi, BN-C307-aroti, BN-C307-dolano, BN-R307-review-300-307

\chapter{Drums and Incense}
\label{ch:bn-drums-and-incense}
\hlchaptermodality{307}

\begin{chapteropening}
\textbf{By the end of this chapter:} \emph{I can describe the drums, gongs and incense of a puja evening.}

Complete the last lesson of chapter 307: I can describe the drums, gongs and incense of a puja evening.
\end{chapteropening}

\section[\texorpdfstring{\bn{ঢাক} (ḍhāk)}{ঢাক (ḍhāk)}]{\bn{ঢাক} (ḍhāk) — a dhak; a big festival drum}

\label{lesson:BN-C307-dhak}



\begin{quote}
Before the new one: say the Bengali for a gift, then the Bengali for to entertain guests.
\end{quote}

\subsection*{You'll want to know: \bn{ঢাক}}
\textbf{\bn{ঢাক}} — \emph{ḍhāk} — \textquotedblleft{}a dhak; a big festival drum\textquotedblright{}.

\begin{quote}
\bn{পুজোর} \bn{সকালে} \bn{একটা} \bn{লোক} \bn{ঢাক} \bn{বাজাচ্ছিল।}
\end{quote}

\begin{quote}
\emph{pujor sôkāle ekṭā lok ḍhāk bājāchchhilô.}
\end{quote}

\begin{quote}
On the morning of the puja a man was playing the dhak.
\end{quote}

\begin{cousinweb}
\textbf{\bn{ঢাক}} (\emph{ḍhāk}) is usually linked with Sanskrit \emph{ḍhakkā}, 'a large drum' (Monier-Williams). Its \bn{ঢ} is the breathy partner of \bn{ড}.
\end{cousinweb}

\begin{grammarlens}[title={What you've built}]
One more word for telling what happened.
\end{grammarlens}

\subsection*{Your turn}
\begin{itemize}
\raggedright
  \item \emph{Say it:} \emph{ḍhāk}
  \item \emph{Say it:} \emph{ḍhāk}, once more
  \item \emph{Say it:} \emph{pujor sôkāle ekṭā lok ḍhāk bājāchchhilô.} — On the morning of the puja a man was playing the dhak.
  \item \emph{Recall:} say the Bengali for a gift, then the Bengali for to entertain guests, then say \emph{ḍhāk} again
\end{itemize}

\begin{tcolorbox}[breakable,colback=teal!4,colframe=teal!35!black,title={Before you move on}]
What does \bn{ঢাক} mean? (\textquotedblleft{}A dhak; a big festival drum\textquotedblright{}.) Say it once more.
\end{tcolorbox}
\section[\texorpdfstring{\bn{কাঁসর} (kā̃sôr)}{কাঁসর (kā̃sôr)}]{\bn{কাঁসর} (kā̃sôr) — a bell-metal gong}

\label{lesson:BN-C307-kasor}



\begin{quote}
Before the new one: say the Bengali for to entertain guests, then the Bengali for a big festival drum.
\end{quote}

\subsection*{You'll want to know: \bn{কাঁসর}}
\textbf{\bn{কাঁসর}} — \emph{kā̃sôr} — \textquotedblleft{}a bell-metal gong\textquotedblright{}.

\begin{quote}
\bn{ঢাকের} \bn{সঙ্গে} \bn{একটা} \bn{ছেলে} \bn{কাঁসর} \bn{বাজাল।}
\end{quote}

\begin{quote}
\emph{ḍhāker sôngge ekṭā chhele kā̃sôr bājālô.}
\end{quote}

\begin{quote}
A boy played the gong along with the dhak.
\end{quote}

\begin{cousinweb}
\textbf{\bn{কাঁসর}} (\emph{kā̃sôr}) is built on \bn{কাঁসা}, bell-metal — Sanskrit \emph{kāṁsya}, 'made of bell-metal' (CDIAL 2987, which lists Bengali \emph{kā̃sā}) — and is a flat gong of it, struck beside the drums.
\end{cousinweb}

\begin{grammarlens}[title={What you've built}]
One more word for telling what happened.
\end{grammarlens}

\subsection*{Your turn}
\begin{itemize}
\raggedright
  \item \emph{Say it:} \emph{kā̃sôr}
  \item \emph{Say it:} \emph{kā̃sôr}, once more
  \item \emph{Say it:} \emph{ḍhāker sôngge ekṭā chhele kā̃sôr bājālô.} — A boy played the gong along with the dhak.
  \item \emph{Recall:} say the Bengali for to entertain guests, then the Bengali for a big festival drum, then say \emph{kā̃sôr} again
\end{itemize}

\begin{tcolorbox}[breakable,colback=teal!4,colframe=teal!35!black,title={Before you move on}]
What does \bn{কাঁসর} mean? (\textquotedblleft{}A bell-metal gong\textquotedblright{}.) Say it once more.
\end{tcolorbox}
\section[\texorpdfstring{\bn{ধুনুচি} (dhunuchi)}{ধুনুচি (dhunuchi)}]{\bn{ধুনুচি} (dhunuchi) — an incense burner}

\label{lesson:BN-C307-dhunuchi}



\begin{quote}
Before the new one: say the Bengali for a big festival drum, then the Bengali for a bell-metal gong.
\end{quote}

\subsection*{You'll want to know: \bn{ধুনুচি}}
\textbf{\bn{ধুনুচি}} — \emph{dhunuchi} — \textquotedblleft{}an incense burner\textquotedblright{}.

\begin{quote}
\bn{পুজোর} \bn{সন্ধ্যায়} \bn{পুজারি} \bn{ধুনুচি} \bn{হাতে} \bn{নিয়ে} \bn{নাচলেন।}
\end{quote}

\begin{quote}
\emph{pujor sôndhyāy pujāri dhunuchi hāte niye nāchlen.}
\end{quote}

\begin{quote}
On the evening of the puja the priest danced with an incense burner in his hand.
\end{quote}

\begin{cousinweb}
\textbf{\bn{ধুনুচি}} (\emph{dhunuchi}) is built on \bn{ধুনো}, the resin burned as incense — Sanskrit \emph{dhūṇa}, 'gum, resin' (CDIAL 6838, which lists Bengali \emph{dhunā}) — and is the clay cup it smoulders in.
\end{cousinweb}

\begin{grammarlens}[title={What you've built}]
One more word for telling what happened.
\end{grammarlens}

\subsection*{Your turn}
\begin{itemize}
\raggedright
  \item \emph{Say it:} \emph{dhunuchi}
  \item \emph{Say it:} \emph{dhunuchi}, once more
  \item \emph{Say it:} \emph{pujor sôndhyāy pujāri dhunuchi hāte niye nāchlen.} — On the evening of the puja the priest danced with an incense burner in his hand.
  \item \emph{Recall:} say the Bengali for a big festival drum, then the Bengali for a bell-metal gong, then say \emph{dhunuchi} again
\end{itemize}

\begin{tcolorbox}[breakable,colback=teal!4,colframe=teal!35!black,title={Before you move on}]
What does \bn{ধুনুচি} mean? (\textquotedblleft{}An incense burner\textquotedblright{}.) Say it once more.
\end{tcolorbox}
\section[\texorpdfstring{\bn{আরতি} (ārôti)}{আরতি (ārôti)}]{\bn{আরতি} (ārôti) — an evening worship with lamps}

\label{lesson:BN-C307-aroti}



\begin{quote}
Before the new one: say the Bengali for a bell-metal gong, then the Bengali for an incense burner.
\end{quote}

\subsection*{You'll want to know: \bn{আরতি}}
\textbf{\bn{আরতি}} — \emph{ārôti} — \textquotedblleft{}an evening worship with lamps\textquotedblright{}.

\begin{quote}
\bn{সন্ধ্যায়} \bn{মন্দিরে} \bn{আরতি} \bn{হয়।}
\end{quote}

\begin{quote}
\emph{sôndhyāy mondire ārôti hôy.}
\end{quote}

\begin{quote}
There is an evening worship with lamps at the temple.
\end{quote}

\begin{cousinweb}
\textbf{\bn{আরতি}} (\emph{ārôti}) comes down from Sanskrit \emph{ārātrika}, the light waved at night before an image, and the rite itself (Monier-Williams; CDIAL 1315, which lists Bengali \emph{ārati}, the ceremony).
\end{cousinweb}

\begin{grammarlens}[title={What you've built}]
One more word for telling what happened.
\end{grammarlens}

\subsection*{Your turn}
\begin{itemize}
\raggedright
  \item \emph{Say it:} \emph{ārôti}
  \item \emph{Say it:} \emph{ārôti}, once more
  \item \emph{Say it:} \emph{sôndhyāy mondire ārôti hôy.} — There is an evening worship with lamps at the temple.
  \item \emph{Recall:} say the Bengali for a bell-metal gong, then the Bengali for an incense burner, then say \emph{ārôti} again
\end{itemize}

\begin{tcolorbox}[breakable,colback=teal!4,colframe=teal!35!black,title={Before you move on}]
What does \bn{আরতি} mean? (\textquotedblleft{}An evening worship with lamps\textquotedblright{}.) Say it once more.
\end{tcolorbox}
\section[\texorpdfstring{\bn{দোলানো} (dolāno)}{দোলানো (dolāno)}]{\bn{দোলানো} (dolāno) — to swing; to sway}

\label{lesson:BN-C307-dolano}



\begin{quote}
Before the new one: say the Bengali for an incense burner, then the Bengali for an evening worship with lamps.
\end{quote}

\subsection*{You'll want to know: \bn{দোলানো}}
\textbf{\bn{দোলানো}} — \emph{dolāno} — \textquotedblleft{}to swing; to sway\textquotedblright{}.

\begin{quote}
\bn{আরতির} \bn{সময়} \bn{পুজারি} \bn{ধুনুচি} \bn{দোলালেন।}
\end{quote}

\begin{quote}
\emph{ārôtir sômôy pujāri dhunuchi dolālen.}
\end{quote}

\begin{quote}
During the evening worship the priest swung the incense burner.
\end{quote}

Past (he/she …): \bn{তিনি} \bn{প্রদীপটা} \bn{দোলালেন}, \emph{tini prôdīpṭā dolālen}.

\begin{cousinweb}
\textbf{\bn{দোলানো}} (\emph{dolāno}) is filed in CDIAL 6453 under a Middle Indo-Aryan \emph{dulati}, 'swings', which lists Bengali \emph{dulā}; \bn{দোলানো} is the -ānō verb, to make something swing.
\end{cousinweb}

\begin{grammarlens}[title={What you've built}]
One more word for telling what happened.
\end{grammarlens}

\subsection*{Your turn}
\begin{itemize}
\raggedright
  \item \emph{Say it:} \emph{dolāno}
  \item \emph{Say it:} \emph{dolāno}, once more
  \item \emph{Say it:} \emph{ārôtir sômôy pujāri dhunuchi dolālen.} — During the evening worship the priest swung the incense burner.
  \item \emph{Recall:} say the Bengali for an incense burner, then the Bengali for an evening worship with lamps, then say \emph{dolāno} again
\end{itemize}

\begin{tcolorbox}[breakable,colback=teal!4,colframe=teal!35!black,title={Before you move on}]
What does \bn{দোলানো} mean? (\textquotedblleft{}To swing; to sway\textquotedblright{}.) Say it once more.
\end{tcolorbox}
\section[(dialogue)]{Looking back: from Readers and Listeners to Drums and Incense}

\label{lesson:BN-R307-review-300-307}



\begin{quote}
Say the Bengali for a reader and a listener.
\end{quote}

\subsection*{Your turn}
Say these aloud:

\begin{itemize}
\raggedright
  \item a reader, a listener, an official notice, to publicize, an issue
  \item an eyewitness, an account, a video, detailed, to verify
  \item a tale, to tell a story, to describe, to bring about, gradually
  \item that day, the past, to go past, over time, to commemorate
  \item a festival, a puja, an idol, a pandal, to celebrate
  \item an oil lamp, fireworks, a floor painting, to glitter, a sky lantern
  \item a guest, to invite, a feast, a gift, to entertain guests
  \item a big festival drum, a bell-metal gong, an incense burner, an evening worship with lamps, to swing
\end{itemize}

\begin{tcolorbox}[breakable,colback=teal!4,colframe=teal!35!black,title={Before you move on}]
A reader? (\textbf{\bn{পাঠক}}, \emph{pāṭhôk}.) A listener? (\textbf{\bn{শ্রোতা}}, \emph{shrotā}.) An official notice? (\textbf{\bn{বিজ্ঞপ্তি}}, \emph{bigyôpti}.) To publicize? (\textbf{\bn{প্রচার} \bn{করা}}, \emph{prôchār kôrā}.) An issue? (\textbf{\bn{সংখ্যা}}, \emph{sôngkhyā}.) An eyewitness? (\textbf{\bn{প্রত্যক্ষদর্শী}}, \emph{prôtyôkkhôdôrshī}.) An account? (\textbf{\bn{বিবরণ}}, \emph{bibôrôn}.) A video? (\textbf{\bn{ভিডিও}}, \emph{bhiḍio}.) Detailed? (\textbf{\bn{বিস্তারিত}}, \emph{bistārit}.) To verify? (\textbf{\bn{যাচাই} \bn{করা}}, \emph{jāchāi kôrā}.) A tale? (\textbf{\bn{কাহিনি}}, \emph{kāhini}.) To tell a story? (\textbf{\bn{শোনানো}}, \emph{shonāno}.) To describe? (\textbf{\bn{বর্ণনা} \bn{করা}}, \emph{bôrnônā kôrā}.) To bring about? (\textbf{\bn{ঘটানো}}, \emph{ghôṭāno}.) Gradually? (\textbf{\bn{ক্রমশ}}, \emph{krômôsh}.) That day? (\textbf{\bn{সেদিন}}, \emph{sedin}.) The past? (\textbf{\bn{অতীত}}, \emph{ôtīt}.) To go past? (\textbf{\bn{পেরোনো}}, \emph{perono}.) Over time? (\textbf{\bn{কালক্রমে}}, \emph{kālôkrôme}.) To commemorate? (\textbf{\bn{স্মরণ} \bn{করা}}, \emph{smôrôn kôrā}.) A festival? (\textbf{\bn{উৎসব}}, \emph{utsôb}.) A puja? (\textbf{\bn{পুজো}}, \emph{pujo}.) An idol? (\textbf{\bn{প্রতিমা}}, \emph{prôtimā}.) A pandal? (\textbf{\bn{প্যান্ডেল}}, \emph{pænḍel}.) To celebrate? (\textbf{\bn{উদযাপন} \bn{করা}}, \emph{udjāpôn kôrā}.) An oil lamp? (\textbf{\bn{প্রদীপ}}, \emph{prôdīp}.) Fireworks? (\textbf{\bn{আতশবাজি}}, \emph{ātôshbāji}.) A floor painting? (\textbf{\bn{আলপনা}}, \emph{ālpônā}.) To glitter? (\textbf{\bn{ঝলমল} \bn{করা}}, \emph{jhôlmôl kôrā}.) A sky lantern? (\textbf{\bn{ফানুস}}, \emph{phānus}.) A guest? (\textbf{\bn{অতিথি}}, \emph{ôtithi}.) To invite? (\textbf{\bn{নিমন্ত্রণ} \bn{করা}}, \emph{nimôntrôn kôrā}.) A feast? (\textbf{\bn{ভোজ}}, \emph{bhoj}.) A gift? (\textbf{\bn{উপহার}}, \emph{upôhār}.) To entertain guests? (\textbf{\bn{আপ্যায়ন} \bn{করা}}, \emph{āpyāyôn kôrā}.) A big festival drum? (\textbf{\bn{ঢাক}}, \emph{ḍhāk}.) A bell-metal gong? (\textbf{\bn{কাঁসর}}, \emph{kā̃sôr}.) An incense burner? (\textbf{\bn{ধুনুচি}}, \emph{dhunuchi}.) An evening worship with lamps? (\textbf{\bn{আরতি}}, \emph{ārôti}.) To swing? (\textbf{\bn{দোলানো}}, \emph{dolāno}.)
\end{tcolorbox}
